\section{Energy space, lifts, and normalization}\label{sec:model}

Let $V=(\Z+\tfrac12)^2$ and orient each nearest-neighbor bond in a positive coordinate direction; $\mathcal E$ contains each undirected bond exactly once. Write $D_eu=u(y)-u(x)$ for $e=(x,y)$. Define the closed gradient subspace first:
\[
 \mathcal G=\overline{\{D\varphi:\varphi\in C_c(V;\R)\}}^{\,\ell^2(\mathcal E)}.
\]
The homogeneous energy space and its energy are
\begin{equation}\label{eq:energy-space}
 \inner uv_\X=\sum_e D_euD_ev,\qquad
 \X=\{u:V\to\R:Du\in\mathcal G\}/\R,
 \qquad\E(u)=\sum_e(1-\cos D_eu).
\end{equation}
The quotient identifies vertex representatives differing by a constant.
The map $D:\X\to\mathcal G$ is an isometric isomorphism: each
$g\in\mathcal G$ has zero oriented sum around every finite cycle, since
finite-cycle sums are continuous on $\ell^2(\mathcal E)$ and vanish on
compact gradients. Integrating $g$ along finite lattice paths therefore
reconstructs a vertex representative, unique modulo a constant.
Thus $\X$ is a Hilbert space, equivalently the homogeneous completion of
compactly supported vertex fields represented modulo constants. Finite-path
vertex differences are continuous in its norm. The energy is finite,
nonnegative, and at most $\nrm u_\X^2/2$.

The physically uniform configurations and admissible barriers are
\begin{align}
 \K&=\{u\in\X:D_eu\in2\pi\Z\text{ for every }e\},\label{eq:vacua}\\
 B_b(u)&=\inf_{\substack{\gamma\in C([0,1];\X)\\
                 \gamma(0)=u,\ \gamma(1)\in\K}}
     \max_{t\in[0,1]}\{E_b(\gamma(t))-E_b(u)\}.\label{eq:barrier}
\end{align}
Write $B_0$ when $b=0$. Replacing the endpoint set by the single representative zero can change the stated problem. Each energy is invariant under addition of an element of $\K$.

Every nonzero $k\in\K$ satisfies $\nrm k_\X\ge4\pi$. Indeed, its integer gradient has finite support, and a nonconstant integer field has at least four bonds with nonzero difference on the square lattice. Thus the quotient map $\X\to\X/\K$ is a local isometry and continuous paths lift uniquely after an initial lift is fixed, by subdividing the compact time interval into neighborhoods smaller than half this separation. Consequently \cref{eq:barrier} also describes paths to the physical vacuum in the phase quotient. Convergence on every fixed finite box is a weaker topology: translating a defect to infinity does not supply a strong energy-space path ending at the vacuum.

We use the Riesz gradient $F(u)=\nabla_\X\E(u)$ and its self-adjoint derivative $H(u)$. The basic differential facts will also specify the constants used in the verification.

\begin{lemma}\label{lem:basic}
The energy is smooth on $\X$, and
\begin{align}
 \inner{F(u)}v_\X&=\sum_e\sin(D_eu)D_ev,\label{eq:firstvar}\\
 \inner{H(u)v}w_\X&=\sum_e\cos(D_eu)D_evD_ew.\label{eq:hessian}
\end{align}
The bond evaluation functional has norm $1/\sqrt2$. In particular,
\begin{equation}\label{eq:basic-bounds}
 \nrm{H(u)}\le1,\qquad
 \nrm{H(u)-H(v)}\le\Lambda\nrm{u-v}_\X,
 \qquad \Lambda=1/\sqrt2.
\end{equation}
For $h\in\X$, the gradient and energy Taylor remainders are bounded by
\begin{gather}\label{eq:remainders}
 \nrm{F(u+h)-F(u)-H(u)h}\le\tfrac\Lambda2\nrm h_\X^2,\qquad\qquad
 \nonumber\\
 |\E(u+h)-\E(u)-\inner{F(u)}h_\X
          -\tfrac12\inner{H(u)h}h_\X|
       \le\tfrac\Lambda6\nrm h_\X^3.
\end{gather}
\end{lemma}
\begin{proof}
The Fourier isometry sends a lattice gradient to $d(k)\widehat u(k)$, with $d_j=e^{ik_j}-1$ and $|d|^2=4-2\cos k_1-2\cos k_2$. The squared norm of one bond evaluation is the normalized torus integral of $|d_j|^2/|d|^2$. The two integrals agree by coordinate exchange, and their sum is one, so each is $1/2$. The value at $k=0$ does not affect the integral. This also follows from the projection construction in \cref{sec:green}.

Since $|\sin t|\le|t|$, the first variation defines a bounded functional. Cauchy--Schwarz gives $\nrm{H(u)}\le1$, and
\[
 |\inner{(H(u)-H(v))a}c_\X|
 \le\nrm{D(u-v)}_\infty\nrm a_\X\nrm c_\X
 \le\Lambda\nrm{u-v}_\X\nrm a_\X\nrm c_\X.
\]
The same estimate bounds the trilinear third derivative by $\Lambda$. For every $j\ge2$, placing two gradient factors in $\ell^2$ and the other $j-2$ factors in $\ell^\infty$ bounds the $j$th derivative by $\Lambda^{j-2}\prod_{i=1}^j\nrm{v_i}_\X$; the same coefficient-continuity argument proves smoothness. Differentiation follows by dominated remainder estimates on the bond sums; continuity follows from uniform convergence of the coefficient increments in $\ell^\infty$. Integration along a line gives \cref{eq:remainders}.
\end{proof}

At the vacuum, $H(0)=I_\X$. A coercivity estimate here is measured against $\nrm{Dv}_2^2$, not against the vertex $\ell^2$ norm. The local vertex Laplacian has spectrum touching zero on the infinite lattice, so this energy-space gap is not an $\ell^2$ spectral gap.

For a plaquette $p$, choose principal bond increments in $(-\pi,\pi)$ and take their oriented sum divided by $2\pi$. If no increment is at a branch endpoint, this defines an integer charge $q_p$. A uniform margin from those endpoints and the bond evaluation bound certify charge preservation under energy-norm corrections. The dipolar state in \cref{sec:instance} has finite total energy and exactly two nonzero charges. We use this charge statement only after its global branch and exterior checks, not as an assumed constraint on escape paths.

\subsection{Dissipation metric and dynamical interpretation}

The choice of metric specifies how the phase field dissipates energy. In
this paper the dissipation potential for an instantaneous velocity $v$ is
\[
 \mathcal R(v)=\frac12\sum_e|D_ev|^2
             =\frac12\nrm v_\X^2.
\]
Thus the velocity is the unique minimizer of
$\mathcal R(v)+DE_b(u)[v]$ over $\X$. Its Euler equation is
\begin{equation}\label{eq:flow}
 \dot u=-\nabla_\X E_b(u).
\end{equation}
Its weak form is
\begin{equation}\label{eq:flow-weak}
 \sum_eD_e\dot u\,D_ev
       =-\sum_e(1+b_e)\sin(D_eu)D_ev
       \quad(v\in\X).
\end{equation}
Writing $D^*$ for the adjoint from bond fields to vertex distributions,
this is $D^*D\dot u=-D^*((1+b)\sin Du)$, understood weakly. Along a
solution,
\[
 \frac{\mathrm d}{\mathrm dt}E_b(u(t))
   =-\nrm{\dot u(t)}_\X^2=-\sum_e|D_e\dot u(t)|^2.
\]
The metric therefore measures dissipation through differences of phase
velocities across bonds, rather than through independent vertex velocities.

This interpretation has a circuit realization. In a finite array of
overdamped resistively shunted Josephson junctions with equal shunt
conductances, a bond voltage is proportional to its phase difference
velocity. After nondimensionalizing, the sum of the resistive and
Josephson bond currents has the form
$i_e=D_e\dot u+J_e\sin(D_eu)$, where $J_e=1+b_e$.
Kirchhoff conservation at every free vertex gives $D^*i=0$ in the absence
of injected current and noise, which is the weak equation above when
the boundary phases are held fixed. The homogeneous current-conservation
equation is given in \cite[Eq.~(13)]{Yoon2000}; fluctuating twists require
the additional equations specified there and are not included here.
The corresponding Joule dissipation is proportional to
$\sum_e|D_e\dot u|^2$. On the infinite lattice we use the finite-energy
weak formulation itself, rather than assuming a thermodynamic limit of
the driven or noisy finite-array model. In particular, no bounded
inverse of the vertex Laplacian on $\ell^2$ is asserted.

The barrier theorem and the dynamical corollary answer distinct questions.
The barrier is determined by the energy and the admissible strong-$\X$
path class in \eqref{eq:barrier}; changing the mobility while keeping
those two objects fixed does not change this variational quantity.
The flow above provides a specified dissipative relaxation model and
the local stability assertion in \cref{cor:flow}. An energy barrier alone
does not validate a connecting orbit, a thermal escape rate, or the
thermodynamic Berezinskii--Kosterlitz--Thouless transition. A different
mobility, noise law, or path topology requires its own dynamical analysis.
